\documentclass{vgtc}                          % final (conference style)
%\documentclass[review]{vgtc}                 % review
%\documentclass[widereview]{vgtc}             % wide-spaced review
%\documentclass[preprint]{vgtc}               % preprint
%\documentclass[electronic]{vgtc}             % electronic version

\let\ifpdf\relax % needed this for compiling with pdflatex -Lane

%% Uncomment one of the lines above depending on where your paper is
%% in the conference process. ``review'' and ``widereview'' are for review
%% submission, ``preprint'' is for pre-publication, and the final version
%% doesn't use a specific qualifier. Further, ``electronic'' includes
%% hyperreferences for more convenient online viewing.

%% Please use one of the ``review'' options in combination with the
%% assigned online id (see below) ONLY if your paper uses a double blind
%% review process. Some conferences, like IEEE Vis and InfoVis, have NOT
%% in the past.

%% Figures should be in CMYK or Grey scale format, otherwise, colour 
%% shifting may occur during the printing process.

%% These three lines bring in essential packages: ``mathptmx'' for Type 1 
%% typefaces, ``graphicx'' for inclusion of EPS figures. and ``times''
%% for proper handling of the times font family.

\usepackage{mathptmx}
\usepackage{graphicx}
\usepackage{times}

%% We encourage the use of mathptmx for consistent usage of times font
%% throughout the proceedings. However, if you encounter conflicts
%% with other math-related packages, you may want to disable it.

%% If you are submitting a paper to a conference for review with a double
%% blind reviewing process, please replace the value ``0'' below with your
%% OnlineID. Otherwise, you may safely leave it at ``0''.
\onlineid{224}

%% declare the category of your paper, only shown in review mode
\vgtccategory{Research}

%% allow for this line if you want the electronic option to work properly
\vgtcinsertpkg

%% In preprint mode you may define your own headline.
%\preprinttext{To appear in an IEEE VGTC sponsored conference.}

%% Paper title.
\title{Markdown to LaTeX!}

%% Author and Affiliation (multiple authors with single affiliations).
\author{John A. Zoidberg \thanks{e-mail: zoidberg@planetexpress.earth}
\and Philip J. Fry       \thanks{e-mail: fry@planetexpress.earth}
}
\affiliation{\scriptsize Deliveries Department \\ Planet Express, Inc.}

%% A teaser figure can be included as follows, but is not recommended since
%% the space is now taken up by a full width abstract.
%\teaser{
%  \includegraphics[width=1.5in]{sample.eps}
%  \caption{Lookit! Lookit!}
%}

%% Abstract section.
\abstract{
Things were did.
} % end of abstract


%% ACM Computing Classification System (CCS). 
%% See <http://www.acm.org/class/1998/> for details.
%% The ``\CCScat'' command takes four arguments.

\CCScatlist{ 
  \CCScat{H.5.2}{Information Interfaces \& Presentations}{User Interfaces - Graphical User Interfaces (GUI)}{};
  \CCScat{I.3.6}{Methodology and Techniques}{Interaction Techniques}{}
}

%% Copyright space is enabled by default as required by guidelines.
%% It is disabled by the 'review' option or via the following command:
% \nocopyrightspace

%%%%%%%%%%%%%%%%%%%%%%%%%%%%%%%%%%%%%%%%%%%%%%%%%%%%%%%%%%%%%%%%
%%%%%%%%%%%%%%%%%%%%%% START OF THE PAPER %%%%%%%%%%%%%%%%%%%%%%
%%%%%%%%%%%%%%%%%%%%%%%%%%%%%%%%%%%%%%%%%%%%%%%%%%%%%%%%%%%%%%%%%
\begin{document}

%% The ``\maketitle'' command must be the first command after the
%% ``\begin{document}'' command. It prepares and prints the title block.

%% the only exception to this rule is the \firstsection command
\firstsection{Introduction}
\maketitle

%Intro body through conclusion comes in here.
We did things.

\begin{figure}[htbp]
\centering
\includegraphics{figs/codementum.png}
\caption{Get to writing!}
\end{figure}

\section{Conclusion}

They were good things\cite{codementum}.



\bibliographystyle{abbrv}
%%use following if all content of bibtex file should be shown
%\nocite{*}
\bibliography{paper}
\end{document}


% Figure Example
\begin{figure}
  \centering
  \includegraphics[width=.7\linewidth]{system.eps}
  \caption{An overview of the visualization system participants used}
  \label{fig:system}
\end{figure}

% Table example (with some small optimizations)
\begin{small}
\begin{table}
    \caption{Results}
    \label{tab:results}
    \scriptsize
    \begin{center}
       \begin{tabular}{| l | c | c | c | c |}
           \hline
           Metric & First & Second\\ \hline
           Accuracy  & 67.71\%  & 67.56\% \\ \hline
           Error  & 32.29\%  & 32.44\% \\ \hline
           Sensitivity  & 61.28\%  & 56.89\% \\ \hline
       \end{tabular}
    \end{center}
\end{table}
\end{small}

% List Example
\begin{itemize}
    \item this
    \item that
    \item win 
\end{itemize}
